\documentclass[11pt]{article}
\usepackage[T1]{fontenc}
\usepackage[a4paper,margin=1in]{geometry}
\usepackage{amsmath,amssymb,amsthm}
\usepackage[hidelinks]{hyperref}
\usepackage{microtype}

\newtheorem*{conjecture}{Source Conjecture 1.1}
\newtheorem*{later}{Later Theorem 1.3}

\setlength{\parindent}{0pt}
\setlength{\parskip}{0.6em}

\title{Literature Status: A Strong Spacetime Variant\\
of the Deift Conjecture Is False}
\author{Run \texttt{fa\_banach\_001} --- lane 17}
\date{12 August 2026}

\begin{document}
\maketitle

\begin{center}
\fbox{\parbox{0.9\textwidth}{
\textbf{Status: literature-implied answer (partial variant).}
The separate later paper arXiv:2209.07501 explicitly cites the counterexample
program arXiv:2111.09345 and disproves the strong spacetime/spatial-preservation
form of the Deift conjecture.  It does not fully settle the source's global
temporal almost-periodicity statement.
}}
\end{center}

\section{Source conjecture and program}

Damanik, Luki\'c, Volberg, and Yuditskii,
\emph{The Deift Conjecture: A Program to Construct a Counterexample},
arXiv:2111.09345v1~\cite{source}, state on arXiv PDF page~2:

\begin{conjecture}
For the KdV equation
\[
  \partial_tu-6u\partial_xu+\partial_x^3u=0
\]
with almost-periodic initial data $u(x,0)=V(x)$, the solution evolves almost
periodically in time.
\end{conjecture}

The paper does not claim a counterexample.  It proposes a coordinated program
involving linearization beyond the direct-Cauchy-theorem regime, construction
of a special Widom domain, stability of time frequencies, and an explicit
description of singular characters.  The last description is identified in
the source as the main obstacle.

\section{Separate later result}

Chapouto, Killip, and Vi\c{s}an,
\emph{Bounded solutions of KdV: uniqueness and the loss of almost
periodicity}, arXiv:2209.07501v1~\cite{laterpaper}, explicitly cite
arXiv:2111.09345 as the preceding program and say that they take a different,
simpler route.  They formulate the strong prediction that the KdV solution is
almost periodic in spacetime, which in particular would make every spatial
profile and every temporal trace almost periodic.  On arXiv PDF page~7 they
state:

\begin{later}
There is a bounded solution $q:[-T,T]\times\mathbb R\to\mathbb R$ of KdV
with almost-periodic initial data for which
$x\mapsto q(t_0,x)$ is not almost periodic at some
$t_0\in[-T,T]$.
\end{later}

The authors immediately identify this as a negative answer to the strong
spacetime formulation.  Their unconditional uniqueness theorem also rules out
replacing this bounded solution by a better-behaved development of the same
initial datum.  Thus the later theorem gives a genuine KdV counterexample to
spatial preservation, and hence to joint spacetime almost periodicity, while
bypassing the source's Widom/DCT program.

\section{Why the scope is partial}

The exact temporal statement in arXiv:2111.09345 is stronger in a different
direction: Bohr almost periodicity in $t$ is a global property on
$\mathbb R$.  Theorem~1.3 of the later paper constructs its KdV solution only
on $[-T,T]$.  On arXiv PDF page~8, the authors explicitly state that they do
not know whether this solution exists globally in time.  They do obtain a
temporal non-almost-periodicity statement for the corresponding \emph{Airy}
evolution, but not a global KdV temporal trace.

Consequently, spatial loss at one time disproves the strong spacetime
formulation but does not, by itself, prove that a globally defined KdV trace
$t\mapsto q(t,x)$ fails to be almost periodic.  Nor does the later paper solve
the source's specific Widom-domain Problem~1.6 or its auxiliary conjectures.
This packet therefore records a partial-variant literature answer, not a full
resolution of Conjecture~1.1.

\section{Search evidence}

The four lightweight run indexes contained no prior treatment of either arXiv
id.  A bounded exact-title, author, Deift-conjecture, and KdV-counterexample
search through 12 August 2026 found arXiv:2209.07501 as the decisive later
paper.  Its citation of arXiv:2111.09345 and its own scope language make the
relationship explicit; the remaining temporal limitation follows from its
local time interval and global-existence disclaimer.  No new mathematical
theorem is claimed here.

\begin{thebibliography}{9}
\bibitem{source}
D. Damanik, M. Luki\'c, A. Volberg, and P. Yuditskii,
\emph{The Deift Conjecture: A Program to Construct a Counterexample},
arXiv:2111.09345v1 (2021).

\bibitem{laterpaper}
A. Chapouto, R. Killip, and M. Vi\c{s}an,
\emph{Bounded solutions of KdV: uniqueness and the loss of almost
periodicity}, arXiv:2209.07501v1 (2022).
\end{thebibliography}

\end{document}
